\chapter{Unit 2: Control Flow, Functions, and Problem-Solving}

\section{Section A: Multiple Choice Practice Questions (60 Curated Problems)}

\noindent\textbf{Q1.} What does the else block execute when the if condition is false? \hfill [Tag: Conditional statements (if-else) | Easy]
\begin{enumerate}[label=(\alph*),topsep=2pt,itemsep=1pt,partopsep=0pt,parsep=0pt]
    \item {The function block}
    \item {The loop block}
    \item {The else block}
    \item {The if block}
\end{enumerate}
\vspace{2mm}
\noindent\textbf{Q2.} Which loop is commonly used to repeat through each item in a sequence? \hfill [Tag: Loops (for, while) | Easy]
\begin{enumerate}[label=(\alph*),topsep=2pt,itemsep=1pt,partopsep=0pt,parsep=0pt]
    \item {A def loop}
    \item {An if loop}
    \item {A for loop}
    \item {A try loop}
\end{enumerate}
\vspace{2mm}
\noindent\textbf{Q3.} What is a nested loop? \hfill [Tag: Nested loops | Easy]
\begin{enumerate}[label=(\alph*),topsep=2pt,itemsep=1pt,partopsep=0pt,parsep=0pt]
    \item {A loop after a function}
    \item {A loop that runs once}
    \item {A loop inside another loop}
    \item {A loop with no condition}
\end{enumerate}
\vspace{2mm}
\noindent\textbf{Q4.} What does the break statement do inside a loop? \hfill [Tag: Break and continue | Easy]
\begin{enumerate}[label=(\alph*),topsep=2pt,itemsep=1pt,partopsep=0pt,parsep=0pt]
    \item {Skips one iteration}
    \item {Restarts the loop}
    \item {Defines a new loop}
    \item {Ends the loop}
\end{enumerate}
\vspace{2mm}
\noindent\textbf{Q5.} What value does this code print? count = 0; count += 1; print(count) \hfill [Tag: Counting | Easy]
\begin{enumerate}[label=(\alph*),topsep=2pt,itemsep=1pt,partopsep=0pt,parsep=0pt]
    \item {2}
    \item {-1}
    \item {0}
    \item {1}
\end{enumerate}
\vspace{2mm}
\noindent\textbf{Q6.} What numbers are produced by range(3) ? \hfill [Tag: range() function | Easy]
\begin{enumerate}[label=(\alph*),topsep=2pt,itemsep=1pt,partopsep=0pt,parsep=0pt]
    \item {0, 1, 2}
    \item {3, 2, 1}
    \item {1, 2, 3}
    \item {0, 1, 2, 3}
\end{enumerate}
\vspace{2mm}
\noindent\textbf{Q7.} Which keyword is used to define a function in Python? \hfill [Tag: Function definitions | Easy]
\begin{enumerate}[label=(\alph*),topsep=2pt,itemsep=1pt,partopsep=0pt,parsep=0pt]
    \item {def}
    \item {define}
    \item {function}
    \item {func}
\end{enumerate}
\vspace{2mm}
\noindent\textbf{Q8.} In greet("Sam") , what is "Sam" called? \hfill [Tag: Arguments | Easy]
\begin{enumerate}[label=(\alph*),topsep=2pt,itemsep=1pt,partopsep=0pt,parsep=0pt]
    \item {A return value}
    \item {A loop counter}
    \item {A condition}
    \item {An argument}
\end{enumerate}
\vspace{2mm}
\noindent\textbf{Q9.} Which keyword sends a value back from a function? \hfill [Tag: Return values | Easy]
\begin{enumerate}[label=(\alph*),topsep=2pt,itemsep=1pt,partopsep=0pt,parsep=0pt]
    \item {give}
    \item {return}
    \item {output}
    \item {send}
\end{enumerate}
\vspace{2mm}
\noindent\textbf{Q10.} What is recursion in programming? \hfill [Tag: Recursion basics | Easy]
\begin{enumerate}[label=(\alph*),topsep=2pt,itemsep=1pt,partopsep=0pt,parsep=0pt]
    \item {A string joining words}
    \item {A loop changing variables}
    \item {A condition using two values}
    \item {A function calling itself}
\end{enumerate}
\vspace{2mm}
\noindent\textbf{Q11.} Which keyword creates a small anonymous function in Python? \hfill [Tag: Lambda function | Easy]
\begin{enumerate}[label=(\alph*),topsep=2pt,itemsep=1pt,partopsep=0pt,parsep=0pt]
    \item {lambda}
    \item {anonymous}
    \item {quick}
    \item {small}
\end{enumerate}
\vspace{2mm}
\noindent\textbf{Q12.} What is the result of "Python"[0:2] ? \hfill [Tag: String slicing | Easy]
\begin{enumerate}[label=(\alph*),topsep=2pt,itemsep=1pt,partopsep=0pt,parsep=0pt]
    \item {Py}
    \item {Python}
    \item {Pyt}
    \item {ho}
\end{enumerate}
\vspace{2mm}
\noindent\textbf{Q13.} Which syntax is an f-string in Python? \hfill [Tag: Advanced string formatting | Easy]
\begin{enumerate}[label=(\alph*),topsep=2pt,itemsep=1pt,partopsep=0pt,parsep=0pt]
    \item {format "Hello \{name\}"}
    \item {string("Hello \{name\}")}
    \item {f"Hello \{name\}"}
    \item {"Hello \{name\}"}
\end{enumerate}
\vspace{2mm}
\noindent\textbf{Q14.} What does the string method .upper() do? \hfill [Tag: String methods | Easy]
\begin{enumerate}[label=(\alph*),topsep=2pt,itemsep=1pt,partopsep=0pt,parsep=0pt]
    \item {Counts the characters}
    \item {Removes spaces}
    \item {Makes letters uppercase}
    \item {Reverses the string}
\end{enumerate}
\vspace{2mm}
\noindent\textbf{Q15.} What does "a,b,c".split(",") return? \hfill [Tag: Splitting and joining strings | Easy]
\begin{enumerate}[label=(\alph*),topsep=2pt,itemsep=1pt,partopsep=0pt,parsep=0pt]
    \item {The string "abc"}
    \item {The list ["a", "b", "c"]}
    \item {The tuple ("a", "b", "c")}
    \item {The number 3}
\end{enumerate}
\vspace{2mm}
\noindent\textbf{Q16.} What is the result of "Hello, \{\}".format("Mia") ? \hfill [Tag: String format method | Easy]
\begin{enumerate}[label=(\alph*),topsep=2pt,itemsep=1pt,partopsep=0pt,parsep=0pt]
    \item {Hello, format}
    \item {Mia, Hello}
    \item {Hello, Mia}
    \item {Hello, \{Mia\}}
\end{enumerate}
\vspace{2mm}
\noindent\textbf{Q17.} What is a regular expression mainly used for? \hfill [Tag: Regular expressions | Easy]
\begin{enumerate}[label=(\alph*),topsep=2pt,itemsep=1pt,partopsep=0pt,parsep=0pt]
    \item {Repeating Python loops}
    \item {Creating numerical lists}
    \item {Matching text patterns}
    \item {Drawing program windows}
\end{enumerate}
\vspace{2mm}
\noindent\textbf{Q18.} What is an algorithm? \hfill [Tag: Algorithmic thinking | Easy]
\begin{enumerate}[label=(\alph*),topsep=2pt,itemsep=1pt,partopsep=0pt,parsep=0pt]
    \item {A type of string}
    \item {A Python comment}
    \item {A named variable}
    \item {A step-by-step procedure}
\end{enumerate}
\vspace{2mm}
\noindent\textbf{Q19.} Which pattern is commonly used to find the total of numbers in a list? \hfill [Tag: Writing Python code for common problem-solving patterns | Easy]
\begin{enumerate}[label=(\alph*),topsep=2pt,itemsep=1pt,partopsep=0pt,parsep=0pt]
    \item {Use a string slice}
    \item {Use a regular expression}
    \item {Use an accumulator}
    \item {Use a nested function definition}
\end{enumerate}
\vspace{2mm}
\noindent\textbf{Q20.} What does the continue statement do inside a loop? \hfill [Tag: Break and continue | Easy]
\begin{enumerate}[label=(\alph*),topsep=2pt,itemsep=1pt,partopsep=0pt,parsep=0pt]
    \item {Skips to the next iteration}
    \item {Repeats the previous iteration forever}
    \item {Exits the current function}
    \item {Stops the program}
\end{enumerate}
\vspace{2mm}
\noindent\textbf{Q21.} What is printed by this code? \hfill [Tag: Conditional statements (if-else) | Medium]
\begin{enumerate}[label=(\alph*),topsep=2pt,itemsep=1pt,partopsep=0pt,parsep=0pt]
    \item {C}
    \item {Nothing}
    \item {B}
    \item {A}
\end{enumerate}
\vspace{2mm}
\noindent\textbf{Q22.} What is the final value of total ? \hfill [Tag: Loops (for, while) | Medium]
\begin{enumerate}[label=(\alph*),topsep=2pt,itemsep=1pt,partopsep=0pt,parsep=0pt]
    \item {8}
    \item {6}
    \item {10}
    \item {12}
\end{enumerate}
\vspace{2mm}
\noindent\textbf{Q23.} How many times is count += 1 executed? \hfill [Tag: Nested loops | Medium]
\begin{enumerate}[label=(\alph*),topsep=2pt,itemsep=1pt,partopsep=0pt,parsep=0pt]
    \item {5}
    \item {6}
    \item {9}
    \item {3}
\end{enumerate}
\vspace{2mm}
\noindent\textbf{Q24.} What is printed by this code? \hfill [Tag: Break and continue | Medium]
\begin{enumerate}[label=(\alph*),topsep=2pt,itemsep=1pt,partopsep=0pt,parsep=0pt]
    \item {1 2 4}
    \item {1 2}
    \item {1 2 4 5}
    \item {1 2 3 4}
\end{enumerate}
\vspace{2mm}
\noindent\textbf{Q25.} What is the value of count after execution? \hfill [Tag: Counting | Medium]
\begin{enumerate}[label=(\alph*),topsep=2pt,itemsep=1pt,partopsep=0pt,parsep=0pt]
    \item {2}
    \item {1}
    \item {4}
    \item {3}
\end{enumerate}
\vspace{2mm}
\noindent\textbf{Q26.} What list is produced by list(range(2, 10, 3)) ? \hfill [Tag: range() function | Medium]
\begin{enumerate}[label=(\alph*),topsep=2pt,itemsep=1pt,partopsep=0pt,parsep=0pt]
    \item {[3, 6, 9]}
    \item {[2, 5, 8, 10]}
    \item {[2, 4, 6, 8]}
    \item {[2, 5, 8]}
\end{enumerate}
\vspace{2mm}
\noindent\textbf{Q27.} What does the following function return when called as square(5) ? \hfill [Tag: Function definitions | Medium]
\begin{enumerate}[label=(\alph*),topsep=2pt,itemsep=1pt,partopsep=0pt,parsep=0pt]
    \item {20}
    \item {55}
    \item {25}
    \item {10}
\end{enumerate}
\vspace{2mm}
\noindent\textbf{Q28.} What is printed by this code? \hfill [Tag: Arguments | Medium]
\begin{enumerate}[label=(\alph*),topsep=2pt,itemsep=1pt,partopsep=0pt,parsep=0pt]
    \item {Mina None and Leo 18}
    \item {Mina 21 and Leo None}
    \item {Mina 18 and Leo 21}
    \item {Mina 21 and Leo 18}
\end{enumerate}
\vspace{2mm}
\noindent\textbf{Q29.} What is printed by this code? \hfill [Tag: Return values | Medium]
\begin{enumerate}[label=(\alph*),topsep=2pt,itemsep=1pt,partopsep=0pt,parsep=0pt]
    \item {None}
    \item {110}
    \item {An error message}
    \item {100}
\end{enumerate}
\vspace{2mm}
\noindent\textbf{Q30.} What is returned by sum\_down(4) ? \hfill [Tag: Recursion basics | Medium]
\begin{enumerate}[label=(\alph*),topsep=2pt,itemsep=1pt,partopsep=0pt,parsep=0pt]
    \item {10}
    \item {24}
    \item {4}
    \item {6}
\end{enumerate}
\vspace{2mm}
\noindent\textbf{Q31.} What is the result of this expression? \hfill [Tag: Lambda function | Medium]
\begin{enumerate}[label=(\alph*),topsep=2pt,itemsep=1pt,partopsep=0pt,parsep=0pt]
    \item {[2, 3, 4, 5]}
    \item {[2, 4, 6, 8]}
    \item {[3, 5, 7, 9]}
    \item {[1, 3, 5, 7]}
\end{enumerate}
\vspace{2mm}
\noindent\textbf{Q32.} What is the value of result ? \hfill [Tag: String slicing | Medium]
\begin{enumerate}[label=(\alph*),topsep=2pt,itemsep=1pt,partopsep=0pt,parsep=0pt]
    \item {"YT"}
    \item {"YH"}
    \item {"YTH"}
    \item {"PTO"}
\end{enumerate}
\vspace{2mm}
\noindent\textbf{Q33.} What is printed by this code? \hfill [Tag: Advanced string formatting | Medium]
\begin{enumerate}[label=(\alph*),topsep=2pt,itemsep=1pt,partopsep=0pt,parsep=0pt]
    \item {Ava: 87.5}
    \item {Ava: 87.46}
    \item {Ava: 87.4}
    \item {Ava: 87}
\end{enumerate}
\vspace{2mm}
\noindent\textbf{Q34.} What is the value of cleaned ? \hfill [Tag: String methods | Medium]
\begin{enumerate}[label=(\alph*),topsep=2pt,itemsep=1pt,partopsep=0pt,parsep=0pt]
    \item {"Python-Programming"}
    \item {"-python-programming-"}
    \item {"python programming"}
    \item {"python-programming"}
\end{enumerate}
\vspace{2mm}
\noindent\textbf{Q35.} What is the result of this code? \hfill [Tag: Splitting and joining strings | Medium]
\begin{enumerate}[label=(\alph*),topsep=2pt,itemsep=1pt,partopsep=0pt,parsep=0pt]
    \item {"red | green | blue"}
    \item {"red,green,blue"}
    \item {"red green blue"}
    \item {"red|green|blue|"}
\end{enumerate}
\vspace{2mm}
\noindent\textbf{Q36.} What is printed by this code? \hfill [Tag: String format method | Medium]
\begin{enumerate}[label=(\alph*),topsep=2pt,itemsep=1pt,partopsep=0pt,parsep=0pt]
    \item {\{\} costs \$\{:.2f\}}
    \item {book costs \$12.50}
    \item {book costs \$12.5}
    \item {book costs \$12}
\end{enumerate}
\vspace{2mm}
\noindent\textbf{Q37.} Which pattern matches a string containing exactly three digits? \hfill [Tag: Regular expressions | Medium]
\begin{enumerate}[label=(\alph*),topsep=2pt,itemsep=1pt,partopsep=0pt,parsep=0pt]
    \item {r"\textasciicircum{}\textbackslash{}\textbackslash{}d\{3\}\$"}
    \item {r"\textbackslash{}\textbackslash{}d+"}
    \item {r"\textbackslash{}\textbackslash{}w\{3\}"}
    \item {r"[0-9]"}
\end{enumerate}
\vspace{2mm}
\noindent\textbf{Q38.} Which strategy correctly finds the largest value in a nonempty list without using max() ? \hfill [Tag: Algorithmic thinking | Medium]
\begin{enumerate}[label=(\alph*),topsep=2pt,itemsep=1pt,partopsep=0pt,parsep=0pt]
    \item {Start with the first value and replace it when a value is larger}
    \item {Start with the last value and replace it when a value is equal}
    \item {Start with 0 and replace it when a value is smaller}
    \item {Sort the list and always select its first value}
\end{enumerate}
\vspace{2mm}
\noindent\textbf{Q39.} Which code correctly counts the number of even values in numbers ? \hfill [Tag: Writing Python code for common problem-solving patterns | Medium]
\begin{enumerate}[label=(\alph*),topsep=2pt,itemsep=1pt,partopsep=0pt,parsep=0pt]
    \item {PYTHON Copy count = 0 for n in numbers: if n \% 2 != 0: count -= 1}
    \item {PYTHON Copy count = 0 for n in numbers: if n \% 2 == 0: count += 1}
    \item {PYTHON Copy count = 1 for n in numbers: if n / 2 == 0: count += 1}
    \item {PYTHON Copy count = 0 for n in numbers: if n \% 2 == 1: count += 1}
\end{enumerate}
\vspace{2mm}
\noindent\textbf{Q40.} Which code correctly creates a list containing the squares of the positive values in numbers ? \hfill [Tag: Writing Python code for common problem-solving patterns | Medium]
\begin{enumerate}[label=(\alph*),topsep=2pt,itemsep=1pt,partopsep=0pt,parsep=0pt]
    \item {PYTHON Copy squares = [n * n for n in numbers if n $>$ 0]}
    \item {PYTHON Copy squares = [n * n for n in numbers if n $<$ 0]}
    \item {PYTHON Copy squares = [n ** 2 for n in numbers if n == 0]}
    \item {PYTHON Copy squares = [n + n for n in numbers if n $>$ 0]}
\end{enumerate}
\vspace{2mm}
\noindent\textbf{Q41.} What is printed by this code? x, y = 8, 3 if x \% y == 0: result = "A" elif x // y == 2 and x $>$ y: result = "B" else: result = "C" print(result) \hfill [Tag: Conditional statements (if-else) | Hard]
\begin{enumerate}[label=(\alph*),topsep=2pt,itemsep=1pt,partopsep=0pt,parsep=0pt]
    \item {B}
    \item {C}
    \item {It raises a ZeroDivisionError}
    \item {A}
\end{enumerate}
\vspace{2mm}
\noindent\textbf{Q42.} What list is produced by list(range(10, -2, -3)) ? \hfill [Tag: range() function | Hard]
\begin{enumerate}[label=(\alph*),topsep=2pt,itemsep=1pt,partopsep=0pt,parsep=0pt]
    \item {[7, 4, 1, -2]}
    \item {[10, 8, 6, 4, 2, 0]}
    \item {[10, 7, 4, 1]}
    \item {[10, 7, 4, 1, -2]}
\end{enumerate}
\vspace{2mm}
\noindent\textbf{Q43.} What is printed by this code? total = 0 for i in range(1, 4): for j in range(i, 4): total += i * j print(total) \hfill [Tag: Nested loops | Hard]
\begin{enumerate}[label=(\alph*),topsep=2pt,itemsep=1pt,partopsep=0pt,parsep=0pt]
    \item {36}
    \item {25}
    \item {30}
    \item {35}
\end{enumerate}
\vspace{2mm}
\noindent\textbf{Q44.} What is printed? values = [] for n in range(2, 10): if n \% 2 == 0: continue if n $>$ 6: break values.append(n) print(values) \hfill [Tag: Break and continue | Hard]
\begin{enumerate}[label=(\alph*),topsep=2pt,itemsep=1pt,partopsep=0pt,parsep=0pt]
    \item {[3, 5, 7]}
    \item {[2, 3, 4, 5, 6]}
    \item {[3, 5]}
    \item {[2, 4, 6]}
\end{enumerate}
\vspace{2mm}
\noindent\textbf{Q45.} Which expression correctly counts the number of overlapping occurrences of "ana" in "banana" ? \hfill [Tag: Counting | Hard]
\begin{enumerate}[label=(\alph*),topsep=2pt,itemsep=1pt,partopsep=0pt,parsep=0pt]
    \item {sum("ana" in "banana"[i:] for i in range(len("banana")))}
    \item {len("banana".split("ana")) - 1}
    \item {"banana".count("ana")}
    \item {sum("banana"[i:i+3] == "ana" for i in range(len("banana") - 2))}
\end{enumerate}
\vspace{2mm}
\noindent\textbf{Q46.} What is printed by this code? def update(items=[]): items.append(len(items)) return items print(update()) print(update()) \hfill [Tag: Function definitions | Hard]
\begin{enumerate}[label=(\alph*),topsep=2pt,itemsep=1pt,partopsep=0pt,parsep=0pt]
    \item {A TypeError occurs because the default list is immutable}
    \item {[0] followed by [0]}
    \item {[0] followed by [1]}
    \item {[1] followed by [1]}
\end{enumerate}
\vspace{2mm}
\noindent\textbf{Q47.} What happens when this code runs? def f(a, b=2, *args, **kwargs): return a + b + len(args) + len(kwargs) print(f(5, 3, 7, 8, x=1)) \hfill [Tag: Arguments | Hard]
\begin{enumerate}[label=(\alph*),topsep=2pt,itemsep=1pt,partopsep=0pt,parsep=0pt]
    \item {It prints 13}
    \item {It prints 10}
    \item {It raises a TypeError because positional arguments follow a default argument}
    \item {It prints 12}
\end{enumerate}
\vspace{2mm}
\noindent\textbf{Q48.} What is printed by this code? def transform(x): if x $<$ 0: return return x * 2 print(transform(-1), transform(4)) \hfill [Tag: Return values | Hard]
\begin{enumerate}[label=(\alph*),topsep=2pt,itemsep=1pt,partopsep=0pt,parsep=0pt]
    \item {It raises a ValueError for the negative argument}
    \item {None None}
    \item {0 8}
    \item {None 8}
\end{enumerate}
\vspace{2mm}
\noindent\textbf{Q49.} What is printed by this function call? def f(n): if n $<$= 1: return 1 return n * f(n - 2) print(f(6)) \hfill [Tag: Recursion basics | Hard]
\begin{enumerate}[label=(\alph*),topsep=2pt,itemsep=1pt,partopsep=0pt,parsep=0pt]
    \item {48}
    \item {6}
    \item {8}
    \item {24}
\end{enumerate}
\vspace{2mm}
\noindent\textbf{Q50.} What is the result of sorted([(1, 4), (2, 1), (1, 2)], key=lambda p: (p[0], -p[1])) ? \hfill [Tag: Lambda function | Hard]
\begin{enumerate}[label=(\alph*),topsep=2pt,itemsep=1pt,partopsep=0pt,parsep=0pt]
    \item {[(1, 4), (2, 1), (1, 2)]}
    \item {[(2, 1), (1, 4), (1, 2)]}
    \item {[(1, 4), (1, 2), (2, 1)]}
    \item {[(1, 2), (1, 4), (2, 1)]}
\end{enumerate}
\vspace{2mm}
\noindent\textbf{Q51.} What is printed by s = "algorithm"; print(s[-2:1:-2]) ? \hfill [Tag: String slicing | Hard]
\begin{enumerate}[label=(\alph*),topsep=2pt,itemsep=1pt,partopsep=0pt,parsep=0pt]
    \item {hio}
    \item {mrt}
    \item {mgt}
    \item {mto}
\end{enumerate}
\vspace{2mm}
\noindent\textbf{Q52.} What is printed by x = 12.3456; print(f"\{x:08.2f\}") ? \hfill [Tag: Advanced string formatting | Hard]
\begin{enumerate}[label=(\alph*),topsep=2pt,itemsep=1pt,partopsep=0pt,parsep=0pt]
    \item {12.35}
    \item {12.345600}
    \item {0012.346}
    \item {00012.35}
\end{enumerate}
\vspace{2mm}
\noindent\textbf{Q53.} What is the result of " A B ".strip().replace(" ", "-").lower() ? \hfill [Tag: String methods | Hard]
\begin{enumerate}[label=(\alph*),topsep=2pt,itemsep=1pt,partopsep=0pt,parsep=0pt]
    \item {"A-B"}
    \item {"a-b"}
    \item {" a-b "}
    \item {"a\textbackslash{} b"}
\end{enumerate}
\vspace{2mm}
\noindent\textbf{Q54.} What is printed by text = "a,,b,"; print("|".join(text.split(","))) ? \hfill [Tag: Splitting and joining strings | Hard]
\begin{enumerate}[label=(\alph*),topsep=2pt,itemsep=1pt,partopsep=0pt,parsep=0pt]
    \item {a||b|}
    \item {a,,b,}
    \item {a|b}
    \item {a| |b|}
\end{enumerate}
\vspace{2mm}
\noindent\textbf{Q55.} What is printed by template = "\{1\}-\{0\}:\{1\}"; print(template.format("x", "y")) ? \hfill [Tag: String format method | Hard]
\begin{enumerate}[label=(\alph*),topsep=2pt,itemsep=1pt,partopsep=0pt,parsep=0pt]
    \item {x-y:y}
    \item {y-x:y}
    \item {x-y:x}
    \item {A KeyError is raised because numeric fields are invalid}
\end{enumerate}
\vspace{2mm}
\noindent\textbf{Q56.} Using Python's re.findall , what does re.findall(r"\textbackslash{}b\textbackslash{}w+\textbackslash{}b", "hi, 42 times!") return? \hfill [Tag: Regular expressions | Hard]
\begin{enumerate}[label=(\alph*),topsep=2pt,itemsep=1pt,partopsep=0pt,parsep=0pt]
    \item {['hi,', '42', 'times!']}
    \item {['hi', 'times']}
    \item {['hi', '42', 'times']}
    \item {['h', 'i', '4', '2', 't', 'i', 'm', 'e', 's']}
\end{enumerate}
\vspace{2mm}
\noindent\textbf{Q57.} A sorted list may contain duplicates. Which strategy finds whether two distinct elements sum to a target in O ( n ) time and O ( 1 ) extra space? \hfill [Tag: Algorithmic thinking | Hard]
\begin{enumerate}[label=(\alph*),topsep=2pt,itemsep=1pt,partopsep=0pt,parsep=0pt]
    \item {Binary-search every possible complement}
    \item {Use two pointers at the beginning and end}
    \item {Try every pair with two nested loops}
    \item {Sort the list first, then scan only adjacent values}
\end{enumerate}
\vspace{2mm}
\noindent\textbf{Q58.} Which code correctly returns the first character in s whose frequency is exactly one, or None if no such character exists? \hfill [Tag: Writing Python code for common problem-solving patterns | Hard]
\begin{enumerate}[label=(\alph*),topsep=2pt,itemsep=1pt,partopsep=0pt,parsep=0pt]
    \item {return min((c for c in s if s.count(c) == 1), default=None)}
    \item {return next((c for c in s if s.count(c) $>$ 1), None)}
    \item {return next((c for c in set(s) if s.count(c) == 1), None)}
    \item {return next((c for c in s if s.count(c) == 1), None)}
\end{enumerate}
\vspace{2mm}
\noindent\textbf{Q59.} What is printed by this code? n = 19 steps = 0 while n $>$ 1: n = n // 2 steps += 1 print(n, steps) \hfill [Tag: Loops (for, while) | Hard]
\begin{enumerate}[label=(\alph*),topsep=2pt,itemsep=1pt,partopsep=0pt,parsep=0pt]
    \item {0 4}
    \item {1 4}
    \item {1 5}
    \item {2 4}
\end{enumerate}
\vspace{2mm}
\noindent\textbf{Q60.} What does pal("racecar", 0, 6) return? def pal(s, left, right): if left $>$= right: return True return s[left] == s[right] and pal(s, left + 1, right - 1) \hfill [Tag: Recursion basics | Hard]
\begin{enumerate}[label=(\alph*),topsep=2pt,itemsep=1pt,partopsep=0pt,parsep=0pt]
    \item {True}
    \item {False}
    \item {It raises an IndexError at the recursive base case}
    \item {None}
\end{enumerate}
\vspace{2mm}

\section{Section B: Comprehensive Subjective Problem with Detailed Analytical Rubric}

\begin{subjectiveq}{2}{1}
Implement a robust log parser \texttt{parse\_access\_logs(raw\_log\_lines)} using regular expressions that extracts IP addresses, HTTP request methods (GET, POST, etc.), endpoint URIs, integer status codes, and response sizes. The function must gracefully filter malformed lines without crashing and return an aggregated frequency table of status codes using recursion.
\end{subjectiveq}

\begin{solbox}
\begin{lstlisting}
import re

LOG_PATTERN = re.compile(
    r'^(\d{1,3}(?:\.\d{1,3}){3})\s+-\s+-\s+\[(.*?)\]\s+"([A-Z]+)\s+([^\s]+)\s+HTTP/[0-9.]+"\s+(\d{3})\s+(\d+|-)'
)

def parse_access_logs(lines):
    parsed = []
    for line in lines:
        match = LOG_PATTERN.match(line.strip())
        if match:
            ip, timestamp, method, endpoint, status, size = match.groups()
            parsed.append({
                "ip": ip, "method": method, "endpoint": endpoint,
                "status": int(status), "size": int(size) if size != '-' else 0
            })
    return parsed

def count_status_recursive(records, idx=0, acc=None):
    if acc is None:
        acc = {}
    if idx >= len(records):
        return acc
    code = records[idx]["status"]
    acc[code] = acc.get(code, 0) + 1
    return count_status_recursive(records, idx + 1, acc)
\end{lstlisting}
\textbf{Analytical Verification:}\\
$\bullet$ Pre-compiling the regex via \texttt{re.compile} reduces parsing latency from $\mathcal{O}(M \cdot K)$ overhead to optimized NFA state evaluation.\\
$\bullet$ Recursion base case handles empty records. For deep logs, tail-call or iterative loops avoid \texttt{RecursionError} (default limit: 1000).
\end{solbox}

\begin{examinertip}
$\bullet$ \textbf{Rubric Breakdown (10 Marks Total):}\\
1. Algorithmic Formulation \& Theoretical Correctness: \textbf{3 Marks}\\
2. Production Python Syntax \& Idiomatic Standards (PEP 8): \textbf{3 Marks}\\
3. Boundary Case \& Exception Robustness: \textbf{2 Marks}\\
4. Asymptotic Time \& Auxiliary Space Complexity Derivation: \textbf{2 Marks}
\end{examinertip}
